\section{总结与延伸}

本讲的逻辑可以压缩成一条逐层变难的链。第一，真实 human intent 不只包含字面指令；第二，RLHF 用 demonstrations、comparisons、reward model 与 RL 把部分隐含偏好变成可优化信号；第三，InstructGPT 与早期 ChatGPT 证明这种信号能显著改善可用性，却保留 hallucination、prompt sensitivity 与代表性问题；第四，当模型完成的任务超过人类独立评价能力，普通 RLHF 的信号会失真；第五，AI-assisted evaluation 试图让机器承担阅读、搜索、计算和 critique，人类集中表达偏好；第六，targeted perturbations、critique experiments 与 discriminator--critique gap 只能逐步测量这条路线是否真的扩展监督。

\begin{importantbox}{整堂课最值得带走的五点}
\begin{enumerate}
  \item Alignment 的对象是 human intent，而不是机械 obedience。
  \item Reward model 是有限比较数据的可复用 proxy，同时也是新的被优化攻击面。
  \item Preference improvement、capability、truthfulness 与 safety 必须分开测量。
  \item Scalable oversight 的关键不是让 AI 替人类决定价值，而是把难评价对象转换成可核验的证据。
  \item 所有递归监督、self-critique、debate 与 interpretability 路线最终都会遇到同一问题：可信 ground truth 或 trusted evaluation signal 从哪里来。
\end{enumerate}
\end{importantbox}

\subsection{证据成熟度阶梯}

读者可以用四层阶梯判断一个 alignment claim 的强度。机制层说明方法怎样工作；当前实验层说明它在受控任务上是否有效；迁移层检验它是否适用于更难、更真实、更具对抗性的任务；保证层才讨论面对比 evaluator 更强的系统时是否仍可靠。本讲的 RLHF 结果主要位于前两层，scalable oversight 结果是第二层的早期 proof of concept，outer/inner alignment 与 interpretability 的最终作用仍属于开放研究。

\begin{warningbox}{不要把研究路线写成既成事实}
课堂给出了一个有希望的方向：用 AI 帮助人类评价 AI。它没有证明评价曲线一定能永远追上能力曲线，也没有证明同一模型会诚实暴露自己的全部知识。高质量讲义应保留这种不确定性，因为问题的难点正是我们无法用当前小实验直接验证未来强系统。
\end{warningbox}

\subsection{拓展阅读}

以下资料均在课堂日期前公开，并直接支撑本讲主线：

\begin{itemize}
  \item \href{https://arxiv.org/abs/1706.03741}{Christiano et al., Deep Reinforcement Learning from Human Preferences}：成对轨迹偏好与在线 reward learning。
  \item \href{https://arxiv.org/abs/1811.07871}{Leike et al., Scalable Agent Alignment via Reward Modeling}：reward modeling 与 recursive assistance 的研究议程。
  \item \href{https://arxiv.org/abs/2203.02155}{Ouyang et al., Training Language Models to Follow Instructions with Human Feedback}：SFT、RM、PPO 与 PPO-ptx 的公开实现和评测。
  \item \href{https://arxiv.org/abs/2009.01325}{Stiennon et al., Learning to Summarize with Human Feedback}：语言生成中的偏好优化。
  \item \href{https://arxiv.org/abs/2206.05802}{Saunders et al., Self-critiquing Models for Assisting Human Evaluators}：critique assistance、targeted flaws 与 discriminator--critique gap。
  \item \href{https://arxiv.org/abs/1805.00899}{Irving et al., AI Safety via Debate}：用对话与对抗论证放大人类判断的另一条路线。
\end{itemize}

最终，alignment 不是把一个 ``goodness score'' 接到模型后面，而是设计一条不会因能力增长而失去可信度的反馈链。讲者在最后没有承诺某个组件会解决全部问题；他给出的更可靠研究姿态是：让评价任务可测、让证据可检查、让失败边界显式，并持续追问模型知道什么、说出什么、实际做了什么。

\subsection{课后自检}

读完本讲后，不应只会复述 RLHF 流程，还应能判断一个新的 alignment 方案究竟改善了哪一段反馈链。下面四个问题可以作为阅读论文、设计实验或评审产品声明时的最低检查表。

\begin{knowledgebox}{用四个问题检查新的 alignment claim}
\begin{enumerate}
  \item 训练信号来自 demonstrations、comparisons、verifiable outcomes，还是另一个未经验证的模型判断？
  \item 实验测量的是 preference、真实正确率、calibration、robustness，还是只测一个容易被优化的 proxy？
  \item 当任务难度超过人类独立评价能力时，系统提供了什么可核验的 evidence，而不只是更流畅的解释？
  \item 失败发生时，能否区分模型不知道、知道但说不出、知道却不愿说，以及 evaluator 没有检查到？
\end{enumerate}
\end{knowledgebox}

\end{document}
